\documentclass[11pt,a4paper]{article}

\usepackage[utf8]{inputenc}
\usepackage[T1]{fontenc}
\usepackage[brazil]{babel}
\usepackage[margin=2.3cm]{geometry}
\usepackage{amsmath,amssymb,amsthm}
\usepackage{enumitem}
\usepackage{booktabs}
\usepackage[dvipsnames]{xcolor}
\usepackage[framemethod=default]{mdframed}
\usepackage{hyperref}
\hypersetup{colorlinks=true,linkcolor=NavyBlue,urlcolor=NavyBlue}

% ---------------- Ambientes ----------------
\theoremstyle{definition}
\newtheorem{definicao}{Definição}[section]
\newtheorem{exemplo}{Exemplo}[section]
\newtheorem{observacao}{Observação}[section]
\theoremstyle{plain}
\newtheorem{teorema}{Teorema}[section]
\newtheorem{lema}[teorema]{Lema}
\newtheorem{corolario}[teorema]{Corolário}
\newtheorem{proposicao}[teorema]{Proposição}

\surroundwithmdframed[linecolor=NavyBlue,linewidth=0.8pt,backgroundcolor=NavyBlue!3]{definicao}
\surroundwithmdframed[linecolor=BrickRed,linewidth=0.8pt,backgroundcolor=BrickRed!3]{teorema}
\surroundwithmdframed[linecolor=ForestGreen,linewidth=0.8pt,backgroundcolor=ForestGreen!3]{exemplo}

% ---------------- Notação ----------------
\newcommand{\E}{\mathbb{E}}
\newcommand{\Var}{\operatorname{Var}}
\newcommand{\Cov}{\operatorname{Cov}}
\newcommand{\EQM}{\operatorname{EQM}}
\newcommand{\B}{\operatorname{B}}
\newcommand{\Prob}{\mathbb{P}}
\newcommand{\R}{\mathbb{R}}
\newcommand{\ind}{\mathbf{1}}
\newcommand{\iid}{\overset{\text{i.i.d.}}{\sim}}
\newcommand{\convP}{\xrightarrow{\;\Prob\;}}
\newcommand{\convMQ}{\xrightarrow{\;\text{MQ}\;}}
\newcommand{\bY}{\bar{Y}}
\newcommand{\just}[1]{\tag*{\footnotesize\textcolor{gray}{[#1]}}}

\title{\textbf{CE309 -- Inferência Estatística}\\[4pt]
\large Guia de Estudos: Propriedades de Estimadores e Método dos Momentos}
\author{Material de apoio para a Lista 6}
\date{}

\begin{document}
\maketitle
\tableofcontents
\newpage

%=====================================================================
\section{Preliminares: o problema de estimação}
%=====================================================================

\begin{definicao}[Modelo estatístico e amostra aleatória]
Seja $Y$ uma variável aleatória com função de densidade (ou de probabilidade) $f(y;\theta)$, em que $\theta \in \Theta \subseteq \R^p$ é um parâmetro \emph{desconhecido} e $\Theta$ é o \emph{espaço paramétrico}. Dizemos que $Y_1,\dots,Y_n$ é uma \emph{amostra aleatória} (a.a.) de $Y$ se as $Y_i$ são independentes e identicamente distribuídas (i.i.d.) com a mesma distribuição de $Y$. A densidade conjunta é
\[
f(y_1,\dots,y_n;\theta) = \prod_{i=1}^n f(y_i;\theta).
\]
\end{definicao}

\begin{definicao}[Estatística, estimador e estimativa]
Uma \emph{estatística} é qualquer função $T = T(Y_1,\dots,Y_n)$ da amostra que \textbf{não depende} de $\theta$. Um \emph{estimador} de $\theta$ (ou de uma função $\tau(\theta)$) é uma estatística usada para aproximar $\theta$. O valor numérico $t = T(y_1,\dots,y_n)$ obtido com os dados observados é a \emph{estimativa}.
\end{definicao}

\begin{observacao}
O estimador é uma \textbf{variável aleatória}; a estimativa é um \textbf{número}. Todas as propriedades deste guia (viés, variância, EQM, consistência) são propriedades da \emph{distribuição amostral} do estimador, isto é, de como ele se comportaria em repetidas amostras. Por isso calculamos $\E_\theta(T)$, $\Var_\theta(T)$ etc., sempre sob o valor verdadeiro $\theta$. O índice $\theta$ em $\E_\theta$ lembra que a esperança é calculada usando $f(\cdot\,;\theta)$.
\end{observacao}

Notações usadas ao longo do texto:
\[
\bY = \frac{1}{n}\sum_{i=1}^n Y_i, \qquad
\mu = \E(Y), \qquad \sigma^2 = \Var(Y), \qquad
Y_{(1)} = \min_i Y_i, \qquad Y_{(n)} = \max_i Y_i .
\]

%=====================================================================
\section{Ferramentas de esperança e variância}
%=====================================================================

Quase tudo nesta lista se reduz a calcular esperanças e variâncias de somas. Revisamos as regras \emph{com demonstração}.

\begin{teorema}[Linearidade da esperança]
Para constantes $a_1,\dots,a_n,b$ e variáveis aleatórias $Y_1,\dots,Y_n$ com esperança finita (\textbf{não é preciso independência}):
\[
\E\Big(b + \sum_{i=1}^n a_i Y_i\Big) = b + \sum_{i=1}^n a_i \E(Y_i).
\]
\end{teorema}
\begin{proof}
Segue da linearidade da integral (ou da soma, no caso discreto): $\E(aX+bW) = \int\!\!\int (ax+bw)f(x,w)\,dx\,dw = a\E(X)+b\E(W)$; o caso geral sai por indução.
\end{proof}

\begin{teorema}[Variância de combinação linear]
\[
\Var\Big(\sum_{i=1}^n a_iY_i\Big) = \sum_{i=1}^n a_i^2\Var(Y_i) + 2\sum_{i<j} a_ia_j\Cov(Y_i,Y_j).
\]
Se as $Y_i$ são independentes, $\Cov(Y_i,Y_j)=0$ para $i\neq j$ e
\[
\Var\Big(\sum_{i=1}^n a_iY_i\Big) = \sum_{i=1}^n a_i^2\Var(Y_i).
\]
\end{teorema}
\begin{proof}
Seja $S=\sum_i a_iY_i$, com $\E(S)=\sum_i a_i\mu_i$. Então $S-\E(S)=\sum_i a_i(Y_i-\mu_i)$ e
\[
\Var(S)=\E\Big[\Big(\sum_i a_i(Y_i-\mu_i)\Big)^2\Big]
=\E\Big[\sum_i\sum_j a_ia_j(Y_i-\mu_i)(Y_j-\mu_j)\Big]
=\sum_i\sum_j a_ia_j\Cov(Y_i,Y_j).
\]
Separando os termos $i=j$ (que dão $a_i^2\Var(Y_i)$) dos termos $i\neq j$ (que aparecem em pares) obtemos a fórmula. Independência implica $\E(Y_iY_j)=\E(Y_i)\E(Y_j)$, logo $\Cov(Y_i,Y_j)=0$.
\end{proof}

\begin{observacao}[Duas fórmulas usadas o tempo todo]
\[
\Var(Y)=\E(Y^2)-[\E(Y)]^2
\qquad\Longleftrightarrow\qquad
\E(Y^2)=\Var(Y)+[\E(Y)]^2 .
\]
\[
\Var(aY+b)=a^2\Var(Y).
\]
\end{observacao}

\begin{teorema}[Média amostral]\label{teo:media}
Se $Y_1,\dots,Y_n$ é a.a. de $Y$ com $\E(Y)=\mu$ e $\Var(Y)=\sigma^2<\infty$ (\textbf{qualquer distribuição}), então
\[
\E(\bY)=\mu, \qquad \Var(\bY)=\frac{\sigma^2}{n}.
\]
\end{teorema}
\begin{proof}
Tomando $a_i=1/n$ nos dois teoremas anteriores:
\[
\E(\bY)=\sum_{i=1}^n\frac1n\,\mu=\frac{n\mu}{n}=\mu,
\qquad
\Var(\bY)=\sum_{i=1}^n\frac1{n^2}\,\sigma^2=\frac{n\sigma^2}{n^2}=\frac{\sigma^2}{n}.
\]
\end{proof}

\begin{lema}[Identidade da soma de quadrados]\label{lema:sq}
\[
\sum_{i=1}^n (Y_i-\bY)^2 = \sum_{i=1}^n Y_i^2 - n\bY^2.
\]
\end{lema}
\begin{proof}
Expandindo o quadrado e usando $\sum_i Y_i = n\bY$:
\[
\sum_{i=1}^n (Y_i^2 - 2Y_i\bY + \bY^2)
= \sum_{i=1}^n Y_i^2 - 2\bY\sum_{i=1}^n Y_i + n\bY^2
= \sum_{i=1}^n Y_i^2 - 2n\bY^2 + n\bY^2
= \sum_{i=1}^n Y_i^2 - n\bY^2 .
\]
\end{proof}

\begin{teorema}[Esperança da soma de quadrados centrada]\label{teo:sq}
Sob as hipóteses do Teorema~\ref{teo:media} (qualquer distribuição com variância finita),
\[
\E\Big[\sum_{i=1}^n (Y_i-\bY)^2\Big] = (n-1)\sigma^2 .
\]
\end{teorema}
\begin{proof}
Pelo Lema~\ref{lema:sq} e linearidade,
$\E\big[\sum_i (Y_i-\bY)^2\big] = \sum_i \E(Y_i^2) - n\E(\bY^2)$.
Pela Observação acima:
\[
\E(Y_i^2)=\sigma^2+\mu^2,
\qquad
\E(\bY^2)=\Var(\bY)+[\E(\bY)]^2=\frac{\sigma^2}{n}+\mu^2 .
\]
Logo
\[
\E\Big[\sum_{i=1}^n (Y_i-\bY)^2\Big]
= n(\sigma^2+\mu^2) - n\Big(\frac{\sigma^2}{n}+\mu^2\Big)
= n\sigma^2 + n\mu^2 - \sigma^2 - n\mu^2
= (n-1)\sigma^2 .
\]
\end{proof}

\begin{observacao}[Por que $n-1$?]
Ao trocar $\mu$ (desconhecido) por $\bY$, a soma $\sum(Y_i-\bY)^2$ fica \emph{sistematicamente menor} que $\sum(Y_i-\mu)^2$, pois $\bY$ é o valor de $c$ que minimiza $\sum(Y_i-c)^2$. Os desvios $Y_i-\bY$ satisfazem a restrição $\sum_i(Y_i-\bY)=0$: só $n-1$ deles são ``livres'' (graus de liberdade).
\end{observacao}

\subsection{Resultados sob normalidade}

\begin{teorema}[Teorema de Fisher / Cochran -- enunciado]\label{teo:fisher}
Se $Y_1,\dots,Y_n \iid N(\mu,\sigma^2)$, então:
\begin{enumerate}[label=(\roman*)]
\item $\bY \sim N(\mu,\sigma^2/n)$;
\item $\displaystyle \frac{\sum_{i=1}^n (Y_i-\bY)^2}{\sigma^2} \sim \chi^2_{n-1}$;
\item $\bY$ e $\sum_{i=1}^n (Y_i-\bY)^2$ são independentes.
\end{enumerate}
\end{teorema}

\begin{proposicao}[Momentos da qui-quadrado]\label{prop:chi}
Se $W\sim\chi^2_k$, então $\E(W)=k$ e $\Var(W)=2k$.
\end{proposicao}
\begin{proof}
Use a representação $W=\sum_{j=1}^k Z_j^2$ com $Z_j\iid N(0,1)$. Basta calcular $\E(Z^2)$ e $\Var(Z^2)$.

\textit{Passo 1.} $\E(Z^2)=\Var(Z)+[\E(Z)]^2 = 1+0 = 1$.

\textit{Passo 2.} $\E(Z^4)$ por integração por partes, com $\phi(z)=\frac{1}{\sqrt{2\pi}}e^{-z^2/2}$ e $\phi'(z)=-z\phi(z)$:
\[
\E(Z^4)=\int_{-\infty}^{\infty} z^3\cdot z\phi(z)\,dz
=\Big[-z^3\phi(z)\Big]_{-\infty}^{\infty}+\int_{-\infty}^{\infty} 3z^2\phi(z)\,dz
=0+3\E(Z^2)=3 .
\]

\textit{Passo 3.} $\Var(Z^2)=\E(Z^4)-[\E(Z^2)]^2=3-1=2$.

\textit{Passo 4.} Por linearidade e independência:
$\E(W)=\sum_{j=1}^k 1 = k$ e $\Var(W)=\sum_{j=1}^k 2=2k$.
\end{proof}

\begin{observacao}[Como usar]
Se você sabe que $W = cT \sim \chi^2_k$ para uma constante $c>0$, então $T=W/c$ e
\[
\E(T)=\frac{k}{c}, \qquad \Var(T)=\frac{2k}{c^2}.
\]
Essa é a técnica padrão para obter esperança e variância de estimadores de $\sigma^2$ sob normalidade.
\end{observacao}

\subsection{Estatísticas de ordem: máximo e mínimo}

\begin{teorema}[Distribuição do máximo e do mínimo]\label{teo:ordem}
Seja $Y_1,\dots,Y_n$ a.a. de $Y$ com f.d.a. $F$ e densidade $f$. Então
\[
F_{Y_{(n)}}(y)=[F(y)]^n, \qquad f_{Y_{(n)}}(y)=n[F(y)]^{n-1}f(y);
\]
\[
F_{Y_{(1)}}(y)=1-[1-F(y)]^n, \qquad f_{Y_{(1)}}(y)=n[1-F(y)]^{n-1}f(y).
\]
\end{teorema}
\begin{proof}
O máximo é $\le y$ se, e somente se, \emph{todas} as observações são $\le y$. Por independência:
\[
\Prob(Y_{(n)}\le y)=\Prob(Y_1\le y,\dots,Y_n\le y)=\prod_{i=1}^n\Prob(Y_i\le y)=[F(y)]^n .
\]
O mínimo é $>y$ se, e somente se, todas são $>y$:
\[
\Prob(Y_{(1)}>y)=\prod_{i=1}^n\Prob(Y_i>y)=[1-F(y)]^n .
\]
As densidades saem derivando em $y$ (regra da cadeia).
\end{proof}

%=====================================================================
\section{Viés e não-vício}
%=====================================================================

\begin{definicao}[Viés]
O \emph{viés} (ou vício) de um estimador $T$ de $\theta$ é
\[
\B_\theta(T)=\E_\theta(T)-\theta .
\]
$T$ é \emph{não-viciado} (não-viesado) para $\theta$ se $\B_\theta(T)=0$, isto é, $\E_\theta(T)=\theta$, \textbf{para todo} $\theta\in\Theta$. Analogamente, $T$ é não-viciado para $\tau(\theta)$ se $\E_\theta(T)=\tau(\theta)$ para todo $\theta$.
\end{definicao}

\begin{definicao}[Assintoticamente não-viciado]
Uma sequência de estimadores $T_n$ é \emph{assintoticamente não-viciada} se $\lim_{n\to\infty}\B_\theta(T_n)=0$ para todo $\theta\in\Theta$.
\end{definicao}

\begin{observacao}[Armadilhas frequentes]
\begin{enumerate}[label=(\alph*)]
\item ``Para todo $\theta$'' é essencial: um estimador constante $T\equiv 5$ tem viés zero quando $\theta=5$, mas não é não-viciado.
\item Não-vício \textbf{não} é preservado por funções não lineares: se $\E(T)=\theta$, em geral $\E[g(T)]\neq g(\theta)$. Exemplo: se $\E(T)=\theta$, então $\E(T^2)=\Var(T)+\theta^2>\theta^2$ sempre que $\Var(T)>0$. Pela desigualdade de Jensen, para $g$ convexa, $\E[g(T)]\ge g(\E T)$.
\item Não-vício é uma propriedade ``em média'': diz que o estimador acerta o alvo na média de infinitas amostras, mas nada diz sobre o quanto ele se espalha. Por isso precisamos também da variância.
\end{enumerate}
\end{observacao}

\begin{proposicao}[Combinações lineares]\label{prop:comblin}
Se $\E(Y_i)=\mu$ para todo $i$, então $\E\big(\sum_i a_iY_i\big)=\mu\sum_i a_i$.
\end{proposicao}
\begin{proof}
Linearidade: $\E\big(\sum_i a_iY_i\big)=\sum_i a_i\E(Y_i)=\sum_i a_i\mu=\mu\sum_i a_i$.
\end{proof}

%=====================================================================
\section{Erro quadrático médio}
%=====================================================================

\begin{definicao}[EQM]
O \emph{erro quadrático médio} de $T$ como estimador de $\theta$ é
\[
\EQM_\theta(T)=\E_\theta\big[(T-\theta)^2\big].
\]
\end{definicao}

\begin{teorema}[Decomposição viés--variância]\label{teo:decomp}
\[
\EQM_\theta(T)=\Var_\theta(T)+[\B_\theta(T)]^2 .
\]
\end{teorema}
\begin{proof}
Somamos e subtraímos $\E_\theta(T)$ dentro do quadrado:
\[
(T-\theta)^2=\big[(T-\E T)+(\E T-\theta)\big]^2
=(T-\E T)^2+2(T-\E T)(\E T-\theta)+(\E T-\theta)^2 .
\]
Tomando esperança: o primeiro termo dá $\Var(T)$; o último é constante e dá $[\B(T)]^2$; o termo cruzado é
\[
2(\E T-\theta)\,\E(T-\E T)=2(\E T-\theta)(\E T-\E T)=0 .
\]
\end{proof}

\begin{corolario}
Se $T$ é não-viciado, $\EQM_\theta(T)=\Var_\theta(T)$.
\end{corolario}

\begin{observacao}[Interpretação]
O EQM mede a distância quadrática média entre o estimador e o alvo. O viés mede erro \emph{sistemático} (acurácia); a variância mede \emph{dispersão} (precisão). O EQM penaliza ambos. Um estimador levemente viciado mas com variância bem menor pode ter EQM menor que um não-viciado (exemplo da Seção~\ref{ex:shrink}).
\end{observacao}

\subsection{Comparação de estimadores}

\begin{definicao}[Comparação via EQM e eficiência relativa]
Sejam $T_1$ e $T_2$ estimadores de $\theta$.
\begin{enumerate}[label=(\roman*)]
\item $T_1$ é \emph{melhor} (uniformemente) que $T_2$ se $\EQM_\theta(T_1)\le\EQM_\theta(T_2)$ para todo $\theta\in\Theta$, com desigualdade estrita para algum $\theta$.
\item A \emph{eficiência relativa} de $T_1$ em relação a $T_2$ é
\[
\operatorname{ef}_\theta(T_1,T_2)=\frac{\EQM_\theta(T_2)}{\EQM_\theta(T_1)} .
\]
Se $\operatorname{ef}>1$, $T_1$ é mais eficiente naquele $\theta$.
\end{enumerate}
\end{definicao}

\begin{observacao}[As curvas podem se cruzar]
Muitas vezes $\EQM_\theta(T_1)<\EQM_\theta(T_2)$ para alguns valores de $\theta$ e o contrário para outros. Nesse caso \textbf{nenhum dos dois é uniformemente melhor}, e a resposta correta numa prova é: ``depende do valor de $\theta$'', identificando a região onde cada um vence. Para encontrá-la, resolva a desigualdade $\EQM(T_1)<\EQM(T_2)$ em $\theta$.
\end{observacao}

\subsection{Limite inferior de Cramér--Rao (ponte para verossimilhança)}

Esta seção dá um padrão de comparação: a menor variância possível para estimadores não-viciados (sob regularidade).

\begin{definicao}[Verossimilhança, escore e informação]
Para uma a.a. com densidade $f(y;\theta)$, $\theta\in\R$:
\[
L(\theta)=\prod_{i=1}^n f(y_i;\theta),\qquad
\ell(\theta)=\log L(\theta)=\sum_{i=1}^n\log f(y_i;\theta),
\]
\[
U(\theta)=\frac{\partial \ell(\theta)}{\partial\theta}\quad\text{(função escore)},\qquad
I_o(\theta)=-\frac{\partial^2 \ell(\theta)}{\partial\theta^2}\quad\text{(informação observada)},
\]
\[
I_e(\theta)=\E_\theta\big[U(\theta)^2\big]=\E_\theta\big[I_o(\theta)\big]\quad\text{(informação esperada de Fisher)}.
\]
\end{definicao}

\begin{teorema}[Desigualdade de Cramér--Rao]
Sob condições de regularidade (em particular, o \textbf{suporte de $f$ não depende de $\theta$} e é permitido derivar sob o sinal de integral), se $T$ é não-viciado para $\tau(\theta)$, então
\[
\Var_\theta(T)\ \ge\ \frac{[\tau'(\theta)]^2}{I_e(\theta)} .
\]
\end{teorema}
\begin{proof}[Esboço]
(i) $\E_\theta[U(\theta)]=\int \frac{\partial}{\partial\theta}f\,dy=\frac{\partial}{\partial\theta}\int f\,dy=0$, logo $\Var(U)=I_e(\theta)$.
(ii) $\Cov(T,U)=\E(TU)=\int T\,\frac{\partial f}{\partial\theta}\,dy=\frac{\partial}{\partial\theta}\E(T)=\tau'(\theta)$.
(iii) Pela desigualdade de Cauchy--Schwarz, $[\Cov(T,U)]^2\le\Var(T)\Var(U)$, isto é, $[\tau'(\theta)]^2\le\Var(T)\,I_e(\theta)$.
\end{proof}

\begin{observacao}
Um estimador não-viciado cuja variância atinge o limite é chamado \emph{eficiente}. Em modelos como $U(0,\theta)$ o suporte depende de $\theta$, a regularidade falha e o limite \textbf{não se aplica} (estimadores podem ter variância de ordem $1/n^2$, menor que a ``usual'' $1/n$).
\end{observacao}

%=====================================================================
\section{Consistência}
%=====================================================================

Agora olhamos para uma \emph{sequência} de estimadores $T_n=T_n(Y_1,\dots,Y_n)$, $n=1,2,\dots$, e perguntamos: com mais dados, $T_n$ se aproxima de $\theta$?

\begin{definicao}[Consistência em média quadrática]
$T_n$ é \emph{médio quadrático consistente} (consistente em MQ) para $\theta$ se
\[
\lim_{n\to\infty}\EQM_\theta(T_n)=\lim_{n\to\infty}\E_\theta\big[(T_n-\theta)^2\big]=0,\qquad\forall\,\theta\in\Theta .
\]
Notação: $T_n\convMQ\theta$.
\end{definicao}

\begin{definicao}[Consistência em probabilidade (fraca)]
$T_n$ é \emph{consistente} (em probabilidade) para $\theta$ se, para todo $\varepsilon>0$,
\[
\lim_{n\to\infty}\Prob_\theta\big(|T_n-\theta|\ge\varepsilon\big)=0
\quad\Longleftrightarrow\quad
\lim_{n\to\infty}\Prob_\theta\big(|T_n-\theta|<\varepsilon\big)=1,
\qquad\forall\,\theta\in\Theta .
\]
Notação: $T_n\convP\theta$.
\end{definicao}

\begin{observacao}[Leitura da definição]
Fixe uma ``margem de erro'' $\varepsilon$ tão pequena quanto quiser. Consistência diz que a probabilidade de o estimador errar por mais que $\varepsilon$ vai a zero quando $n$ cresce. Repare que é uma afirmação sobre \textbf{probabilidades}, não sobre cada realização.
\end{observacao}

\subsection{As desigualdades que fazem o trabalho}

\begin{teorema}[Desigualdade de Markov]
Se $X\ge 0$ e $a>0$, então $\Prob(X\ge a)\le \dfrac{\E(X)}{a}$.
\end{teorema}
\begin{proof}
Como $X\ge 0$, vale pontualmente $X\ge a\,\ind\{X\ge a\}$ (se $X\ge a$, o lado direito é $a\le X$; caso contrário é $0\le X$). Tomando esperança: $\E(X)\ge a\,\E[\ind\{X\ge a\}]=a\,\Prob(X\ge a)$.
\end{proof}

\begin{corolario}[Chebyshev generalizada]\label{cor:cheb}
Para qualquer $c\in\R$ e $\varepsilon>0$:
\[
\Prob(|T-c|\ge\varepsilon)=\Prob\big((T-c)^2\ge\varepsilon^2\big)\le\frac{\E[(T-c)^2]}{\varepsilon^2}.
\]
Com $c=\theta$: $\ \Prob_\theta(|T-\theta|\ge\varepsilon)\le\EQM_\theta(T)/\varepsilon^2$.
Com $c=\E(T)$: $\ \Prob(|T-\E T|\ge\varepsilon)\le\Var(T)/\varepsilon^2$ (Chebyshev clássica).
\end{corolario}
\begin{proof}
Aplique Markov a $X=(T-c)^2\ge 0$ com $a=\varepsilon^2$; os eventos $\{|T-c|\ge\varepsilon\}$ e $\{(T-c)^2\ge\varepsilon^2\}$ são iguais.
\end{proof}

\subsection{Relações entre os tipos de consistência}

\begin{teorema}[MQ $\Rightarrow$ probabilidade]\label{teo:mq_p}
Se $T_n\convMQ\theta$, então $T_n\convP\theta$.
\end{teorema}
\begin{proof}
Pelo Corolário~\ref{cor:cheb}, para todo $\varepsilon>0$:
\[
0\le\Prob_\theta(|T_n-\theta|\ge\varepsilon)\le\frac{\EQM_\theta(T_n)}{\varepsilon^2}\xrightarrow[n\to\infty]{}0 .
\]
Pelo teorema do confronto, a probabilidade tende a zero.
\end{proof}

\begin{corolario}[Critério prático]\label{cor:criterio}
$\EQM_\theta(T_n)\to 0$ se, e somente se, $\B_\theta(T_n)\to 0$ \textbf{e} $\Var_\theta(T_n)\to 0$.
Logo: se $T_n$ é (assintoticamente) não-viciado e $\Var(T_n)\to 0$, então $T_n$ é consistente em MQ e, portanto, em probabilidade.
\end{corolario}
\begin{proof}
Pela decomposição, $\EQM=\Var+\B^2$ é soma de dois termos não negativos; a soma vai a zero se, e somente se, cada parcela vai a zero.
\end{proof}

\begin{observacao}[A recíproca é falsa]
Consistência em probabilidade \textbf{não} implica consistência em MQ. Contraexemplo: seja
\[
\Prob(T_n=\theta+n)=\tfrac1n,\qquad \Prob(T_n=\theta)=1-\tfrac1n .
\]
Para $0<\varepsilon<1$ e $n\ge 2$: $\Prob(|T_n-\theta|\ge\varepsilon)=\Prob(T_n=\theta+n)=1/n\to0$, logo $T_n\convP\theta$. Mas
\[
\EQM(T_n)=\E[(T_n-\theta)^2]=n^2\cdot\tfrac1n+0\cdot\big(1-\tfrac1n\big)=n\to\infty .
\]
\end{observacao}

\subsection{Lei dos Grandes Números e funções contínuas}

\begin{teorema}[Lei Fraca dos Grandes Números]\label{teo:lgn}
Se $Y_1,\dots,Y_n$ é a.a. com $\E(Y)=\mu$ e $\Var(Y)=\sigma^2<\infty$, então $\bY\convP\mu$. Mais geralmente, se $\E(Y^{2k})<\infty$, o $k$-ésimo momento amostral $M_k=\frac1n\sum_i Y_i^k$ satisfaz $M_k\convP\E(Y^k)$.
\end{teorema}
\begin{proof}
Pelo Teorema~\ref{teo:media}, $\E(\bY)=\mu$ e $\Var(\bY)=\sigma^2/n$. Por Chebyshev,
$\Prob(|\bY-\mu|\ge\varepsilon)\le\dfrac{\sigma^2}{n\varepsilon^2}\to 0$.
Para $M_k$, aplique o mesmo argumento às variáveis $X_i=Y_i^k$, que são i.i.d. com média $\E(Y^k)$ e variância $\E(Y^{2k})-[\E(Y^k)]^2<\infty$.
\end{proof}

\begin{teorema}[Aplicação contínua]\label{teo:cont}
Se $T_n\convP\theta$ e $g$ é contínua em $\theta$, então $g(T_n)\convP g(\theta)$. O resultado vale também para vetores: se $(T_{1n},\dots,T_{pn})\convP(\theta_1,\dots,\theta_p)$ e $g:\R^p\to\R$ é contínua nesse ponto, então $g(T_{1n},\dots,T_{pn})\convP g(\theta_1,\dots,\theta_p)$.
\end{teorema}
\begin{proof}[Ideia]
Continuidade em $\theta$: dado $\varepsilon>0$ existe $\delta>0$ tal que $|t-\theta|<\delta\Rightarrow|g(t)-g(\theta)|<\varepsilon$. Logo
$\{|g(T_n)-g(\theta)|\ge\varepsilon\}\subseteq\{|T_n-\theta|\ge\delta\}$ e
$\Prob(|g(T_n)-g(\theta)|\ge\varepsilon)\le\Prob(|T_n-\theta|\ge\delta)\to 0$.
\end{proof}

\begin{observacao}
Consequências úteis: se $T_n\convP\theta$ e $a_n\to a$ (sequência numérica), então $a_nT_n\convP a\theta$; somas e produtos de sequências consistentes são consistentes para a soma e o produto dos limites.
\end{observacao}

\subsection{Quando um estimador \emph{não} é consistente}

\begin{proposicao}[Estimador que não usa a informação crescente]\label{prop:naocons}
Seja $T$ um estimador cuja distribuição \textbf{não depende de $n$} (por exemplo, uma função de um número fixo de observações). Se $T$ não é degenerado em $\theta$, isto é, existe $\varepsilon_0>0$ com $p_0=\Prob_\theta(|T-\theta|\ge\varepsilon_0)>0$, então $T$ não é consistente.
\end{proposicao}
\begin{proof}
Para todo $n$, $\Prob_\theta(|T_n-\theta|\ge\varepsilon_0)=p_0$, uma constante positiva. Logo o limite é $p_0\neq 0$, violando a definição (que exige limite zero para \emph{todo} $\varepsilon$).
\end{proof}

\begin{observacao}
Para provar \textbf{não}-consistência basta exibir \emph{um} $\varepsilon>0$ para o qual $\Prob(|T_n-\theta|\ge\varepsilon)\not\to0$. Mostrar apenas que o EQM não vai a zero \textbf{não basta}, pois MQ é mais forte que probabilidade (pela observação da recíproca falsa).
\end{observacao}

\subsection{Consistência via f.d.a. (útil para máximos e mínimos)}

Quando $T_n$ é um máximo ou mínimo, muitas vezes é mais fácil calcular $\Prob(|T_n-\theta|\ge\varepsilon)$ \emph{diretamente} pela f.d.a. (Teorema~\ref{teo:ordem}) do que passar pelo EQM. O roteiro é:
\begin{enumerate}[label=\arabic*.]
\item Observe se $T_n$ fica sempre de um lado de $\theta$ (por exemplo $T_n\le\theta$); isso simplifica o valor absoluto.
\item Escreva $\Prob(|T_n-\theta|\ge\varepsilon)$ em termos de $F_{T_n}$.
\item Use que $r^n\to0$ quando $0\le r<1$.
\end{enumerate}
Veja o Exemplo 2 (Seção~\ref{ex:expo}), item (d), onde isso é feito para um mínimo.

%=====================================================================
\section{Método dos momentos}
%=====================================================================

\begin{definicao}[Momentos populacionais e amostrais]
O $k$-ésimo momento populacional (em torno da origem) é $\mu_k'(\theta)=\E_\theta(Y^k)$. O $k$-ésimo momento amostral é
\[
M_k=\frac1n\sum_{i=1}^n Y_i^k .
\]
Em particular $M_1=\bY$.
\end{definicao}

\begin{definicao}[Estimador de momentos]
Se $\theta=(\theta_1,\dots,\theta_p)$, o \emph{estimador de momentos} $\tilde\theta$ é a solução (em $\theta$) do sistema
\[
\mu_1'(\theta)=M_1,\qquad \mu_2'(\theta)=M_2,\qquad \dots,\qquad \mu_p'(\theta)=M_p .
\]
Usa-se a notação $\tilde\theta$ para distinguir do EMV $\hat\theta$.
\end{definicao}

\begin{mdframed}[linecolor=Gray,backgroundcolor=Gray!5]
\textbf{Roteiro de prova -- método dos momentos}
\begin{enumerate}[label=\arabic*.]
\item Conte quantos parâmetros desconhecidos há ($p$).
\item Calcule os $p$ primeiros momentos populacionais $\E(Y),\dots,\E(Y^p)$ como funções de $\theta$, \textbf{fazendo as integrais} (ou consultando a tabela da Seção~\ref{sec:tabela}). Se $\E(Y)$ não depende de $\theta$ (ou for zero por simetria), ele não serve: passe para o próximo momento.
\item Iguale cada momento populacional ao amostral correspondente.
\item Resolva o sistema para $\theta$ e troque $\theta$ por $\tilde\theta$.
\item Confira se a estimativa pertence a $\Theta$ e é compatível com os dados.
\end{enumerate}
\end{mdframed}

\begin{teorema}[Propriedades gerais]
\begin{enumerate}[label=(\roman*)]
\item Cada $M_k$ é não-viciado para $\mu_k'$: $\E(M_k)=\frac1n\sum_i\E(Y_i^k)=\mu_k'$.
\item Se $\tilde\theta=g(M_1,\dots,M_p)$ com $g$ contínua, então $\tilde\theta$ é consistente para $\theta$ (LGN + aplicação contínua, Teoremas~\ref{teo:lgn} e~\ref{teo:cont}).
\item Em geral, $\tilde\theta$ é \textbf{viciado} quando $g$ é não linear (pela observação sobre não-vício e funções não lineares).
\end{enumerate}
\end{teorema}

\begin{observacao}[Limitações do método]
\begin{enumerate}[label=(\alph*)]
\item \textbf{Não unicidade:} o mesmo parâmetro pode ser estimado igualando momentos diferentes (por exemplo, na Poisson $\E(Y)=\Var(Y)=\lambda$), gerando estimadores distintos.
\item \textbf{Incompatibilidade com os dados:} a estimativa pode cair fora de $\Theta$ ou contradizer a amostra (exemplo da Seção~\ref{ex:unifdesloc}).
\item \textbf{Não usa a verossimilhança:} nem sempre é eficiente; costuma servir como ponto de partida (por exemplo, valor inicial para algoritmos de máxima verossimilhança).
\end{enumerate}
\end{observacao}

%=====================================================================
\section{Tabela de momentos de referência}\label{sec:tabela}
%=====================================================================

\begin{center}
\renewcommand{\arraystretch}{1.5}
\begin{tabular}{@{}llll@{}}
\toprule
Distribuição & Densidade / f.p. & $\E(Y)$ & $\Var(Y)$\\
\midrule
$\text{Ber}(p)$ & $p^y(1-p)^{1-y}$, $y\in\{0,1\}$ & $p$ & $p(1-p)$\\
$\text{Bin}(n,p)$ & $\binom ny p^y(1-p)^{n-y}$ & $np$ & $np(1-p)$\\
$\text{Poi}(\lambda)$ & $e^{-\lambda}\lambda^y/y!$ & $\lambda$ & $\lambda$\\
$N(\mu,\sigma^2)$ & $\frac{1}{\sqrt{2\pi}\sigma}e^{-(y-\mu)^2/(2\sigma^2)}$ & $\mu$ & $\sigma^2$\\
$\text{Exp}(\lambda)$ (taxa) & $\lambda e^{-\lambda y}$, $y>0$ & $1/\lambda$ & $1/\lambda^2$\\
$\text{Gama}(\alpha,\beta)$ (taxa) & $\frac{\beta^\alpha}{\Gamma(\alpha)}y^{\alpha-1}e^{-\beta y}$, $y>0$ & $\alpha/\beta$ & $\alpha/\beta^2$\\
$U(a,b)$ & $\frac{1}{b-a}\ind\{a<y<b\}$ & $(a+b)/2$ & $(b-a)^2/12$\\
$\chi^2_k$ & $\text{Gama}(k/2,1/2)$ & $k$ & $2k$\\
\bottomrule
\end{tabular}
\end{center}

\begin{observacao}[Integral gama -- a ferramenta para momentos de densidades do tipo $y^a e^{-by}$]
Para $\alpha>0$, $\beta>0$:
\[
\int_0^\infty y^{\alpha-1}e^{-\beta y}\,dy=\frac{\Gamma(\alpha)}{\beta^\alpha},
\qquad \Gamma(m)=(m-1)!\ \text{ para } m\in\mathbb N .
\]
Prova: substitua $u=\beta y$, $dy=du/\beta$: $\int_0^\infty (u/\beta)^{\alpha-1}e^{-u}\,du/\beta=\beta^{-\alpha}\Gamma(\alpha)$.
Ao calcular $\E(Y^k)$ de uma densidade com fator $y^a e^{-by}$, junte as potências de $y$ e reconheça uma integral gama.
\end{observacao}

%=====================================================================
\section{Caixa de ferramentas: propriedades numeradas}\label{sec:props}
%=====================================================================

Nos exemplos da próxima seção, \textbf{cada passagem algébrica vem com uma etiqueta cinza à direita}, como \textcolor{gray}{[P1]}, indicando qual propriedade abaixo foi usada. A ideia é que você consiga refazer qualquer linha sozinho, sabendo exatamente \emph{por que} ela vale.

\subsection*{Esperança e variância}
\begin{enumerate}[label=\textbf{(P\arabic*)},leftmargin=3.2em]
\item \textbf{Linearidade da esperança.} $\E(aX+b)=a\E(X)+b$ e $\E\big(\sum_i X_i\big)=\sum_i\E(X_i)$. Constantes ``saem'' da esperança e a esperança de uma soma é a soma das esperanças (não exige independência).
\item \textbf{Constantes na variância.} $\Var(aX+b)=a^2\Var(X)$. A constante sai \emph{ao quadrado}; somar uma constante $b$ não altera a dispersão.
\item \textbf{Variância da soma de independentes.} Se $X_1,\dots,X_n$ são independentes, $\Var\big(\sum_i X_i\big)=\sum_i\Var(X_i)$ (as covariâncias são nulas).
\item \textbf{Fórmula de atalho.} $\Var(X)=\E(X^2)-[\E(X)]^2$, ou seja, $\E(X^2)=\Var(X)+[\E(X)]^2$.
\item \textbf{Esperança de constante.} $\E(c)=c$ e $\Var(c)=0$.
\item \textbf{Independência.} $\E(XW)=\E(X)\E(W)$ e $\Prob(A\cap B)=\Prob(A)\Prob(B)$ para eventos ligados a variáveis independentes.
\item \textbf{Esperança de uma função (LOTUS).} $\E[g(X)]=\int g(x)f(x)\,dx$ (contínua) ou $\sum_x g(x)p(x)$ (discreta).
\item \textbf{Identicamente distribuídas.} Em uma a.a., todas as $Y_i$ têm a mesma distribuição de $Y$; logo $\E(Y_i)=\E(Y)$, $\Var(Y_i)=\Var(Y)$ e $\Prob(Y_i>y)=\Prob(Y>y)$ para todo $i$.
\end{enumerate}

\subsection*{Somatórios e produtórios}
\begin{enumerate}[label=\textbf{(S\arabic*)},leftmargin=3.2em]
\item $\sum_{i=1}^n c = nc$ \ (somar a mesma constante $n$ vezes).
\item $\sum_i c\,x_i = c\sum_i x_i$ \ (constante sai do somatório).
\item $\sum_i (x_i+w_i)=\sum_i x_i+\sum_i w_i$.
\item $\sum_{i=1}^n x_i = n\bar x$ \ (definição da média: $\bar x=\frac1n\sum_i x_i$).
\item $\prod_{i=1}^n c = c^n$ e $\prod_i (a_ib_i)=\big(\prod_i a_i\big)\big(\prod_i b_i\big)$.
\end{enumerate}

\subsection*{Logaritmo e exponencial (sempre $\log=\ln$)}
\begin{enumerate}[label=\textbf{(L\arabic*)},leftmargin=3.2em]
\item $\log(ab)=\log a+\log b$; portanto $\log\prod_i a_i=\sum_i\log a_i$.
\item $\log(a^b)=b\log a$.
\item $\log(e^x)=x$; \ $e^ae^b=e^{a+b}$; \ $(e^a)^n=e^{na}$.
\item $\log(a/b)=\log a-\log b$.
\end{enumerate}

\subsection*{Derivadas}
\begin{enumerate}[label=\textbf{(D\arabic*)},leftmargin=3.2em]
\item $\frac{d}{dx}(cx)=c$ e $\frac{d}{dx}c=0$ \ (o que não depende da variável é constante).
\item $\frac{d}{dx}\log x=\frac1x$ \ e \ $\frac{d}{dx}x^{k}=kx^{k-1}$ (em particular $\frac{d}{dx}x^{-1}=-x^{-2}$).
\item Derivada da soma é a soma das derivadas.
\end{enumerate}

\subsection*{Integrais}
\begin{enumerate}[label=\textbf{(I\arabic*)},leftmargin=3.2em]
\item \textbf{Integral gama.} Para $\alpha>0$ e $\beta>0$: $\displaystyle\int_0^\infty y^{\alpha-1}e^{-\beta y}\,dy=\frac{\Gamma(\alpha)}{\beta^\alpha}$.
\item \textbf{Propriedades da função gama.} $\Gamma(\alpha+1)=\alpha\,\Gamma(\alpha)$; \ $\Gamma(1)=1$; \ $\Gamma(m)=(m-1)!$ para $m$ inteiro positivo. Ex.: $\Gamma(2)=1$, $\Gamma(3)=2$, $\Gamma(4)=6$.
\item Constante sai da integral: $\int c\,g(y)\,dy=c\int g(y)\,dy$.
\item \textbf{Polinômios.} $\int_a^b y^k\,dy=\Big[\frac{y^{k+1}}{k+1}\Big]_a^b=\frac{b^{k+1}-a^{k+1}}{k+1}$.
\item \textbf{Exponencial.} $\int_0^y \frac1\beta e^{-t/\beta}\,dt=\Big[-e^{-t/\beta}\Big]_0^y=1-e^{-y/\beta}$.
\end{enumerate}

\subsection*{Probabilidade}
\begin{enumerate}[label=\textbf{(Pr\arabic*)},leftmargin=3.6em]
\item \textbf{Complemento.} $\Prob(A^c)=1-\Prob(A)$; em particular $\Prob(Y>y)=1-F(y)$.
\item \textbf{União disjunta.} Se $A\cap B=\emptyset$, $\Prob(A\cup B)=\Prob(A)+\Prob(B)$.
\item \textbf{Valor absoluto.} $|X-c|\ge\varepsilon \iff X\le c-\varepsilon$ \ ou \ $X\ge c+\varepsilon$ (dois eventos disjuntos).
\item \textbf{FGM.} $M_X(t)=\E(e^{tX})$. A FGM de uma soma de independentes é o produto das FGMs, e a FGM determina a distribuição (unicidade).
\item \textbf{Mínimo e máximo.} $\{Y_{(1)}>y\}=\{Y_1>y,\dots,Y_n>y\}$ \ e \ $\{Y_{(n)}\le y\}=\{Y_1\le y,\dots,Y_n\le y\}$.
\end{enumerate}

\subsection*{Limites e desigualdades}
\begin{enumerate}[label=\textbf{(Li\arabic*)},leftmargin=3.6em]
\item Para $c$ constante (não depende de $n$): $\lim_{n\to\infty} c/n=0$ e $\lim_{n\to\infty}c=c$.
\item \textbf{Confronto.} Se $0\le a_n\le b_n$ e $b_n\to0$, então $a_n\to0$.
\item Multiplicar ou dividir uma desigualdade por número \textbf{positivo} preserva o sentido; somar ou subtrair qualquer número também.
\item Se $a,b\ge0$: $a<b\iff a^2<b^2$, e $\sqrt{a^2}=|a|$.
\end{enumerate}

\subsection*{Resultados do próprio guia que serão citados}
\textbf{[Decomp]} $\EQM=\Var+\B^2$ (Teorema~\ref{teo:decomp}). \quad
\textbf{[Cheb]} $\Prob(|T-\theta|\ge\varepsilon)\le\EQM(T)/\varepsilon^2$ (Corolário~\ref{cor:cheb}). \quad
\textbf{[LGN]} $\bY\convP\E(Y)$ (Teorema~\ref{teo:lgn}). \quad
\textbf{[Cont]} aplicação contínua (Teorema~\ref{teo:cont}). \quad
\textbf{[Lema-SQ]} $\sum_i(Y_i-\bY)^2=\sum_iY_i^2-n\bY^2$ (Lema~\ref{lema:sq}). \quad
\textbf{[Teo-SQ]} $\E\big[\sum_i(Y_i-\bY)^2\big]=(n-1)\sigma^2$ (Teorema~\ref{teo:sq}).

%=====================================================================
\section{Exemplos resolvidos à mão, passo a passo}
%=====================================================================

Os exemplos usam modelos \emph{diferentes} dos da lista, para que as técnicas fiquem claras sem entregar as respostas. Leia cada linha e, antes de olhar a etiqueta, tente dizer qual propriedade justifica a passagem.

%---------------------------------------------------------------
\subsection{Exemplo 1 -- Poisson: as cinco propriedades e Cramér--Rao}\label{ex:poisson}

\textbf{Enunciado.} Seja $Y_1,\dots,Y_n\iid\text{Poi}(\lambda)$, com $f(y;\lambda)=\dfrac{e^{-\lambda}\lambda^y}{y!}$, $y=0,1,2,\dots$, e o estimador $\hat\lambda=\bY=\frac1n\sum_{i=1}^nY_i$.

\subsubsection*{Passo 0 -- de onde vêm $\E(Y)=\lambda$ e $\Var(Y)=\lambda$}

Não decore: veja de onde sai. Vamos usar a série de Taylor $\sum_{k=0}^\infty\frac{\lambda^k}{k!}=e^\lambda$.

\textit{Esperança.} O termo $y=0$ vale zero (está multiplicado por $y=0$), então a soma começa em $y=1$. Como $\frac{y}{y!}=\frac{1}{(y-1)!}$ para $y\ge1$:
\begin{align*}
\E(Y)&=\sum_{y=0}^\infty y\,\frac{e^{-\lambda}\lambda^y}{y!} \just{P7}\\
&=\sum_{y=1}^\infty \frac{e^{-\lambda}\lambda^y}{(y-1)!} \just{termo $y=0$ é nulo; $y/y!=1/(y-1)!$}\\
&=\lambda e^{-\lambda}\sum_{y=1}^\infty \frac{\lambda^{y-1}}{(y-1)!} \just{$\lambda^y=\lambda\cdot\lambda^{y-1}$; S2}\\
&=\lambda e^{-\lambda}\sum_{k=0}^\infty \frac{\lambda^{k}}{k!} \just{troca $k=y-1$}\\
&=\lambda e^{-\lambda}e^{\lambda}=\lambda . \just{série de Taylor; L3}
\end{align*}

\textit{Variância.} O truque é calcular $\E[Y(Y-1)]$, porque $\frac{y(y-1)}{y!}=\frac{1}{(y-2)!}$ para $y\ge2$. Repetindo o raciocínio:
\begin{align*}
\E[Y(Y-1)]&=\sum_{y=2}^\infty \frac{e^{-\lambda}\lambda^y}{(y-2)!}
=\lambda^2e^{-\lambda}\sum_{k=0}^\infty\frac{\lambda^k}{k!}=\lambda^2 . \just{mesmo argumento, $k=y-2$}
\end{align*}
Como $\E[Y(Y-1)]=\E(Y^2)-\E(Y)$ (por P1), temos $\E(Y^2)=\lambda^2+\lambda$. Então
\begin{align*}
\Var(Y)&=\E(Y^2)-[\E(Y)]^2 \just{P4}\\
&=\lambda^2+\lambda-\lambda^2=\lambda .
\end{align*}

\subsubsection*{(a) Não-vício}
\textbf{Objetivo:} mostrar que $\E(\hat\lambda)=\lambda$ para todo $\lambda>0$.
\begin{align*}
\E(\hat\lambda)&=\E\Big(\frac1n\sum_{i=1}^nY_i\Big) \just{definição de $\hat\lambda$}\\
&=\frac1n\,\E\Big(\sum_{i=1}^nY_i\Big) \just{P1: constante $1/n$ sai}\\
&=\frac1n\sum_{i=1}^n\E(Y_i) \just{P1: esperança da soma}\\
&=\frac1n\sum_{i=1}^n\lambda \just{P8 + Passo 0}\\
&=\frac1n\cdot n\lambda \just{S1}\\
&=\lambda .
\end{align*}
Logo $\B_\lambda(\hat\lambda)=\E(\hat\lambda)-\lambda=\lambda-\lambda=0$ para todo $\lambda>0$: $\hat\lambda$ é não-viciado.

\subsubsection*{(b) Variância}
\begin{align*}
\Var(\hat\lambda)&=\Var\Big(\frac1n\sum_{i=1}^nY_i\Big)\\
&=\frac1{n^2}\,\Var\Big(\sum_{i=1}^nY_i\Big) \just{P2: $1/n$ sai \textbf{ao quadrado}}\\
&=\frac1{n^2}\sum_{i=1}^n\Var(Y_i) \just{P3: as $Y_i$ são independentes}\\
&=\frac1{n^2}\sum_{i=1}^n\lambda \just{P8 + Passo 0}\\
&=\frac1{n^2}\cdot n\lambda \just{S1}\\
&=\frac{\lambda}{n}. \just{simplifica $n/n^2=1/n$}
\end{align*}
\textbf{Erro comum:} esquecer o quadrado em P2 e escrever $\frac1n\cdot n\lambda=\lambda$. Outro erro: usar P3 sem citar a independência, que é exatamente o que a justifica.

\subsubsection*{(c) Erro quadrático médio}
\begin{align*}
\EQM_\lambda(\hat\lambda)&=\Var(\hat\lambda)+[\B_\lambda(\hat\lambda)]^2 \just{Decomp}\\
&=\frac{\lambda}{n}+0^2 \just{itens (a) e (b)}\\
&=\frac{\lambda}{n}.
\end{align*}

\subsubsection*{(d) Consistência em média quadrática}
Fixamos um valor qualquer de $\lambda>0$. Ele \textbf{não muda com $n$}, então funciona como uma constante no limite:
\begin{align*}
\lim_{n\to\infty}\EQM_\lambda(\hat\lambda)&=\lim_{n\to\infty}\frac{\lambda}{n}=0 . \just{Li1 com $c=\lambda$}
\end{align*}
Como isso vale para todo $\lambda>0$, $\hat\lambda$ é médio quadrático consistente.

\subsubsection*{(e) Consistência em probabilidade}
\textbf{Objetivo:} para todo $\varepsilon>0$, mostrar que $\Prob(|\hat\lambda-\lambda|\ge\varepsilon)\to0$.
\begin{align*}
\Prob(|\hat\lambda-\lambda|\ge\varepsilon)
&\le\frac{\EQM_\lambda(\hat\lambda)}{\varepsilon^2} \just{Cheb}\\
&=\frac{\lambda/n}{\varepsilon^2}=\frac{\lambda}{n\varepsilon^2}. \just{item (c)}
\end{align*}
Com $\lambda$ e $\varepsilon$ fixos, $\frac{\lambda}{\varepsilon^2}$ é uma constante e $\frac{\lambda}{n\varepsilon^2}\to0$ (Li1). Como toda probabilidade é $\ge0$, temos
\[
0\le\Prob(|\hat\lambda-\lambda|\ge\varepsilon)\le\frac{\lambda}{n\varepsilon^2}\to0,
\]
e pelo confronto (Li2) a probabilidade tende a zero. Portanto $\hat\lambda\convP\lambda$.

\textit{Atalho legítimo em prova:} ``como $\hat\lambda$ é consistente em MQ (item d), pelo Teorema~\ref{teo:mq_p} é consistente em probabilidade''. Mas é bom saber escrever a cadeia acima, porque é ela que \emph{prova} o teorema.

\subsubsection*{(f) Verossimilhança, escore, $I_o$, $I_e$ e Cramér--Rao}

\textit{Verossimilhança.} Por independência, a densidade conjunta é o produto:
\[
L(\lambda)=\prod_{i=1}^n\frac{e^{-\lambda}\lambda^{y_i}}{y_i!}.
\]

\textit{Log-verossimilhança.} O log transforma o produto em soma:
\begin{align*}
\ell(\lambda)&=\log\prod_{i=1}^n\frac{e^{-\lambda}\lambda^{y_i}}{y_i!}
=\sum_{i=1}^n\log\Big(\frac{e^{-\lambda}\lambda^{y_i}}{y_i!}\Big) \just{L1}\\
&=\sum_{i=1}^n\Big[\log e^{-\lambda}+\log\lambda^{y_i}-\log(y_i!)\Big] \just{L1 e L4}\\
&=\sum_{i=1}^n\Big[-\lambda+y_i\log\lambda-\log(y_i!)\Big] \just{L3 e L2}\\
&=\sum_{i=1}^n(-\lambda)+\sum_{i=1}^ny_i\log\lambda-\sum_{i=1}^n\log(y_i!) \just{S3}\\
&=-n\lambda+\Big(\sum_{i=1}^ny_i\Big)\log\lambda-\sum_{i=1}^n\log(y_i!). \just{S1 e S2}
\end{align*}

\textit{Escore.} Derivamos termo a termo em relação a $\lambda$ (D3). Os $y_i$ são números observados, logo constantes:
\begin{align*}
\frac{d}{d\lambda}(-n\lambda)&=-n \just{D1}\\
\frac{d}{d\lambda}\Big[\Big(\sum_iy_i\Big)\log\lambda\Big]&=\Big(\sum_iy_i\Big)\cdot\frac1\lambda \just{D1 e D2}\\
\frac{d}{d\lambda}\Big[\sum_i\log(y_i!)\Big]&=0 \just{D1: não depende de $\lambda$}
\end{align*}
\[
U(\lambda)=-n+\frac{\sum_{i=1}^ny_i}{\lambda}.
\]

\textit{Informação observada.} Derivamos de novo e trocamos o sinal. Escreva $\frac{\sum_iy_i}{\lambda}=\big(\sum_iy_i\big)\lambda^{-1}$:
\begin{align*}
\frac{d^2\ell}{d\lambda^2}=\frac{dU}{d\lambda}&=0+\Big(\sum_iy_i\Big)\cdot(-1)\lambda^{-2} \just{D1 e D2}\\
&=-\frac{\sum_iy_i}{\lambda^2},
\end{align*}
\[
I_o(\lambda)=-\frac{d^2\ell}{d\lambda^2}=\frac{\sum_{i=1}^ny_i}{\lambda^2}.
\]

\textit{Informação esperada.} Trocamos os valores observados $y_i$ pelas variáveis aleatórias $Y_i$ e tomamos esperança:
\begin{align*}
I_e(\lambda)&=\E\Big[\frac{\sum_iY_i}{\lambda^2}\Big]
=\frac1{\lambda^2}\sum_{i=1}^n\E(Y_i) \just{P1}\\
&=\frac{n\lambda}{\lambda^2}=\frac{n}{\lambda}. \just{P8, S1}
\end{align*}
\textit{Conferência pela outra definição, $I_e=\E[U^2]$.} Primeiro, $\E[U]=-n+\frac{n\lambda}{\lambda}=0$. Então $\E[U^2]=\Var(U)$ (P4), e
\[
\Var(U)=\Var\Big(-n+\frac1\lambda\sum_iY_i\Big)=\frac1{\lambda^2}\Var\Big(\sum_iY_i\Big)=\frac{n\lambda}{\lambda^2}=\frac{n}{\lambda}
\]
usando P2 (a constante $-n$ some e $1/\lambda$ sai ao quadrado) e P3. As duas definições batem.

\textit{Limite de Cramér--Rao.} Aqui o alvo é o próprio $\lambda$, então $\tau(\lambda)=\lambda$ e $\tau'(\lambda)=1$:
\[
\text{LICR}=\frac{[\tau'(\lambda)]^2}{I_e(\lambda)}=\frac{1}{n/\lambda}=\frac{\lambda}{n}.
\]
Como $\Var(\hat\lambda)=\lambda/n$ (item b) é \emph{igual} ao limite, $\hat\lambda$ é \textbf{eficiente}. Nenhum estimador não-viciado de $\lambda$ tem variância menor.

\textit{Bônus: o EMV.} Igualando o escore a zero:
\[
-n+\frac{\sum_iy_i}{\lambda}=0
\iff\frac{\sum_iy_i}{\lambda}=n
\iff\lambda=\frac{\sum_iy_i}{n}=\bar y .
\]
A segunda derivada $-\sum_iy_i/\lambda^2<0$ confirma que é um máximo. Aqui o EMV coincide com o estimador de momentos ($\E(Y)=\lambda=\bY$), mas isso \emph{nem sempre} acontece.

\subsubsection*{(g) Numérico}
Amostra $(2,\,0,\,3,\,1,\,4)$, com $n=5$:
\begin{align*}
\sum_iy_i&=2+0+3+1+4=10,\\
\hat\lambda&=\frac{10}{5}=2,\\
\widehat{\Var}(\hat\lambda)&=\frac{\hat\lambda}{n}=\frac{2}{5}=0{,}4 \quad\text{(trocamos o $\lambda$ desconhecido pela estimativa)},\\
\text{erro-padrão}&=\sqrt{0{,}4}\approx0{,}632 .
\end{align*}

%---------------------------------------------------------------
\subsection{Exemplo 2 -- Exponencial: dois estimadores não-viciados, só um consistente}\label{ex:expo}

\textbf{Enunciado.} Seja $Y_1,\dots,Y_n\iid$ Exponencial com \emph{média} $\beta$, isto é, $f(y;\beta)=\frac1\beta e^{-y/\beta}$ para $y>0$. Compare $T_1=\bY$ e $T_2=n\,Y_{(1)}$, em que $Y_{(1)}=\min(Y_1,\dots,Y_n)$.

\subsubsection*{Passo 0 -- momentos e f.d.a. de $Y$}

\textit{Esperança.} Reescrevemos o integrando no formato da integral gama $y^{\alpha-1}e^{-(\text{taxa})\,y}$. Aqui $y=y^{2-1}$ e a taxa é $1/\beta$:
\begin{align*}
\E(Y)&=\int_0^\infty y\cdot\frac1\beta e^{-y/\beta}\,dy \just{P7}\\
&=\frac1\beta\int_0^\infty y^{2-1}e^{-(1/\beta)y}\,dy \just{I3}\\
&=\frac1\beta\cdot\frac{\Gamma(2)}{(1/\beta)^2} \just{I1 com $\alpha=2$, taxa $1/\beta$}\\
&=\frac1\beta\cdot\frac{1}{1/\beta^2} \just{I2: $\Gamma(2)=1!=1$}\\
&=\frac1\beta\cdot\beta^2=\beta .
\end{align*}

\textit{Segundo momento.} Agora $y^2=y^{3-1}$:
\begin{align*}
\E(Y^2)&=\frac1\beta\int_0^\infty y^{3-1}e^{-(1/\beta)y}\,dy
=\frac1\beta\cdot\frac{\Gamma(3)}{(1/\beta)^3} \just{P7, I3, I1 com $\alpha=3$}\\
&=\frac1\beta\cdot2\beta^3=2\beta^2 . \just{I2: $\Gamma(3)=2!=2$}
\end{align*}

\textit{Variância.}
\[
\Var(Y)=\E(Y^2)-[\E(Y)]^2=2\beta^2-\beta^2=\beta^2 . \just{P4}
\]

\textit{F.d.a. e sobrevivência.} Para $y>0$:
\begin{align*}
F(y)&=\int_0^y\frac1\beta e^{-t/\beta}\,dt=1-e^{-y/\beta}, \just{I5}\\
\Prob(Y>y)&=1-F(y)=e^{-y/\beta}. \just{Pr1}
\end{align*}

\subsubsection*{(a) A distribuição de $Y_{(1)}$ e de $T_2$}
\textbf{Ideia-chave:} o mínimo é maior que $y$ se, e somente se, \emph{todas} as observações são maiores que $y$.
\begin{align*}
\Prob(Y_{(1)}>y)&=\Prob(Y_1>y,\;Y_2>y,\;\dots,\;Y_n>y) \just{Pr5}\\
&=\prod_{i=1}^n\Prob(Y_i>y) \just{P6: independência}\\
&=\prod_{i=1}^n e^{-y/\beta} \just{P8 + Passo 0}\\
&=\big(e^{-y/\beta}\big)^n \just{S5}\\
&=e^{-ny/\beta}. \just{L3}
\end{align*}
Compare com a sobrevivência de uma exponencial de média $m$, que é $e^{-y/m}$. Escrevendo $e^{-ny/\beta}=e^{-y/(\beta/n)}$, vemos que $Y_{(1)}$ é exponencial com \textbf{média $\beta/n$}.

Agora $T_2=nY_{(1)}$. Para $t>0$:
\begin{align*}
\Prob(T_2>t)&=\Prob(nY_{(1)}>t)\\
&=\Prob\Big(Y_{(1)}>\frac tn\Big) \just{Li3: dividir por $n>0$}\\
&=e^{-n\,(t/n)/\beta} \just{resultado acima com $y=t/n$}\\
&=e^{-t/\beta}. \just{$n\cdot\frac tn=t$}
\end{align*}
Ou seja, \textbf{$T_2$ tem exatamente a mesma distribuição de uma única observação $Y$, qualquer que seja $n$}. Multiplicar por $n$ desfaz o encolhimento do mínimo.

\subsubsection*{(b) Viés e variância de cada estimador}
\textit{$T_1=\bY$.} Repetindo os passos do Exemplo~1 (a) e (b), agora com $\E(Y_i)=\beta$ e $\Var(Y_i)=\beta^2$:
\[
\E(T_1)=\frac1n\sum_{i=1}^n\beta=\beta,
\qquad
\Var(T_1)=\frac1{n^2}\sum_{i=1}^n\beta^2=\frac{n\beta^2}{n^2}=\frac{\beta^2}{n}.
\]
\textit{$T_2=nY_{(1)}$.} Pelo item (a), $T_2$ tem a mesma distribuição de $Y$, então herda os momentos do Passo 0:
\[
\E(T_2)=\beta,\qquad\Var(T_2)=\beta^2 .
\]
Também dá para ver a esperança por P1: $\E(T_2)=n\,\E(Y_{(1)})=n\cdot\frac{\beta}{n}=\beta$.

\textbf{Conclusão:} $\B(T_1)=\B(T_2)=0$, ou seja, \emph{ambos são não-viciados}.

\subsubsection*{(c) EQM e eficiência relativa}
Como ambos são não-viciados, pela Decomp o EQM é só a variância:
\[
\EQM(T_1)=\frac{\beta^2}{n},\qquad\EQM(T_2)=\beta^2 .
\]
\begin{align*}
\operatorname{ef}(T_1,T_2)&=\frac{\EQM(T_2)}{\EQM(T_1)}=\frac{\beta^2}{\beta^2/n} \just{definição}\\
&=\beta^2\cdot\frac{n}{\beta^2} \just{dividir por fração = multiplicar pelo inverso}\\
&=n .
\end{align*}
Para $n\ge2$, $\EQM(T_1)<\EQM(T_2)$ para \textbf{todo} $\beta>0$, então $T_1$ é uniformemente melhor.

\subsubsection*{(d) Consistência}
\textit{$T_1$.} $\EQM(T_1)=\beta^2/n\to0$ (Li1), logo $T_1$ é consistente em MQ e, portanto, em probabilidade.

\textit{$T_2$.} Vamos calcular a probabilidade de erro diretamente. Tome $0<\varepsilon<\beta$ (assim $\beta-\varepsilon>0$):
\begin{align*}
\Prob(|T_2-\beta|\ge\varepsilon)&=\Prob(T_2\le\beta-\varepsilon)+\Prob(T_2\ge\beta+\varepsilon) \just{Pr3 e Pr2}\\
&=F(\beta-\varepsilon)+e^{-(\beta+\varepsilon)/\beta} \just{$T_2$ tem a f.d.a. $F$ do Passo 0}\\
&=\Big(1-e^{-(\beta-\varepsilon)/\beta}\Big)+e^{-(\beta+\varepsilon)/\beta}\\
&=1-e^{-1+\varepsilon/\beta}+e^{-1-\varepsilon/\beta}. \just{$\frac{\beta\pm\varepsilon}{\beta}=1\pm\frac\varepsilon\beta$}
\end{align*}
Repare que \textbf{não aparece $n$} nessa expressão. Ela é estritamente positiva, porque já a segunda parcela $e^{-1-\varepsilon/\beta}$ é positiva. Logo o limite quando $n\to\infty$ é essa mesma constante positiva, e não zero: \textbf{$T_2$ não é consistente}.

\textit{Numérico}, com $\beta=1$ e $\varepsilon=0{,}5$:
\[
1-e^{-0{,}5}+e^{-1{,}5}\approx1-0{,}6065+0{,}2231=0{,}6166 .
\]
Mesmo com um milhão de observações, há cerca de $62\%$ de chance de $T_2$ errar por $0{,}5$ ou mais.

\textbf{Moral:} não-vício não garante consistência. O que importa é se a dispersão do estimador encolhe com $n$; a de $T_2$ não encolhe.

%---------------------------------------------------------------
\subsection{Exemplo 3 -- Um estimador viciado pode ter EQM menor}\label{ex:shrink}

\textbf{Enunciado.} Seja $Y_1,\dots,Y_n\iid N(\mu,\sigma^2)$ com $\sigma^2$ conhecido. Considere $T_c=c\,\bY$ com $0<c\le1$ (``encolher'' a média em direção a zero).

\subsubsection*{Viés}
\begin{align*}
\E(T_c)&=\E(c\bY)=c\,\E(\bY) \just{P1}\\
&=c\mu, \just{Teorema~\ref{teo:media}}\\
\B(T_c)&=\E(T_c)-\mu=c\mu-\mu=(c-1)\mu . \just{coloca $\mu$ em evidência}
\end{align*}
Se $c<1$ e $\mu\neq0$, o estimador é viciado.

\subsubsection*{Variância}
\begin{align*}
\Var(T_c)&=\Var(c\bY)=c^2\Var(\bY) \just{P2}\\
&=c^2\frac{\sigma^2}{n}. \just{Teorema~\ref{teo:media}}
\end{align*}

\subsubsection*{EQM}
\begin{align*}
\EQM_\mu(T_c)&=\Var(T_c)+[\B(T_c)]^2 \just{Decomp}\\
&=\frac{c^2\sigma^2}{n}+[(c-1)\mu]^2\\
&=\frac{c^2\sigma^2}{n}+(c-1)^2\mu^2 . \just{$(ab)^2=a^2b^2$}
\end{align*}
Leia a fórmula: diminuir $c$ \emph{reduz} a primeira parcela (variância) mas \emph{aumenta} a segunda (viés$^2$). É o \textbf{compromisso viés--variância}.

\subsubsection*{Comparação numérica: $n=10$, $\sigma^2=1$, $c=1$ contra $c=0{,}9$}
\begin{align*}
\EQM(T_1)&=\frac{1^2\cdot1}{10}+(1-1)^2\mu^2=0{,}1,\\
\EQM(T_{0,9})&=\frac{(0{,}9)^2\cdot1}{10}+(0{,}9-1)^2\mu^2=\frac{0{,}81}{10}+(-0{,}1)^2\mu^2=0{,}081+0{,}01\mu^2 .
\end{align*}
Quando $T_{0,9}$ é melhor? Resolvemos a desigualdade passo a passo:
\begin{align*}
0{,}081+0{,}01\mu^2&<0{,}1\\
0{,}01\mu^2&<0{,}019 \just{Li3: subtrai $0{,}081$}\\
\mu^2&<1{,}9 \just{Li3: divide por $0{,}01>0$}\\
|\mu|&<\sqrt{1{,}9}\approx1{,}378 . \just{Li4}
\end{align*}
\textbf{Conclusão:} se $|\mu|<1{,}378$, o estimador viciado $T_{0,9}$ tem EQM menor; se $|\mu|>1{,}378$, $\bY$ vence. As curvas se cruzam, então \emph{nenhum é uniformemente melhor}. Como $\mu$ é desconhecido, não sabemos de antemão em qual região estamos.

\subsubsection*{Bônus: o $c$ ótimo (e por que não dá para usá-lo)}
Derivamos o EQM em relação a $c$ e igualamos a zero. Pela regra da cadeia, $\frac{d}{dc}(c-1)^2=2(c-1)$:
\begin{align*}
\frac{d}{dc}\EQM&=\frac{2c\sigma^2}{n}+2(c-1)\mu^2=0\\
c\frac{\sigma^2}{n}+c\mu^2-\mu^2&=0 \just{divide por 2 e distribui}\\
c\Big(\frac{\sigma^2}{n}+\mu^2\Big)&=\mu^2 \just{coloca $c$ em evidência}\\
c^*&=\frac{\mu^2}{\mu^2+\sigma^2/n}.
\end{align*}
O $c^*$ depende do $\mu$ desconhecido, então $c^*\bY$ \textbf{não é uma estatística}. Lembre que estimador não pode depender do parâmetro.

\subsubsection*{Consistência de $T_c$ com $c<1$ fixo}
O viés $(c-1)\mu$ \textbf{não depende de $n$}. Logo $\EQM\to(c-1)^2\mu^2\neq0$ quando $\mu\ne0$, e não há consistência em MQ. Também não há consistência em probabilidade: por LGN e Cont, $T_c=c\bY\convP c\mu\neq\mu$. Ganhar em amostras pequenas não garante bom comportamento assintótico.

%---------------------------------------------------------------
\subsection{Exemplo 4 -- Momentos: exponencial com taxa $\lambda$, incluindo o viés}\label{ex:mmexp}

\textbf{Enunciado.} Seja $Y_1,\dots,Y_n\iid\text{Exp}(\lambda)$, com $f(y;\lambda)=\lambda e^{-\lambda y}$ para $y>0$. Encontre o estimador de momentos de $\lambda$ e estude suas propriedades.

\subsubsection*{Passo 1 -- momento populacional}
Há um parâmetro ($p=1$), então basta o primeiro momento.

\textit{Via integral gama:}
\begin{align*}
\E(Y)&=\int_0^\infty y\cdot\lambda e^{-\lambda y}\,dy \just{P7}\\
&=\lambda\int_0^\infty y^{2-1}e^{-\lambda y}\,dy \just{I3; $y=y^{2-1}$}\\
&=\lambda\cdot\frac{\Gamma(2)}{\lambda^2} \just{I1: $\alpha=2$, taxa $\lambda$}\\
&=\lambda\cdot\frac{1}{\lambda^2}=\frac1\lambda . \just{I2}
\end{align*}

\textit{Via integração por partes (para conferir).} Com $u=y$, $dv=\lambda e^{-\lambda y}dy$, temos $du=dy$ e $v=-e^{-\lambda y}$:
\begin{align*}
\E(Y)&=\Big[-y\,e^{-\lambda y}\Big]_0^\infty+\int_0^\infty e^{-\lambda y}\,dy
=(0-0)+\Big[-\frac1\lambda e^{-\lambda y}\Big]_0^\infty
=0+\Big(0+\frac1\lambda\Big)=\frac1\lambda .
\end{align*}
No primeiro colchete, $y\,e^{-\lambda y}\to0$ quando $y\to\infty$ porque a exponencial decai mais rápido do que $y$ cresce.

\subsubsection*{Passo 2 -- igualar e resolver}
O método manda igualar o momento populacional ao amostral, trocando $\lambda$ por $\tilde\lambda$:
\begin{align*}
\frac{1}{\tilde\lambda}&=\bY\\
1&=\tilde\lambda\,\bY \just{multiplica os dois lados por $\tilde\lambda$}\\
\tilde\lambda&=\frac{1}{\bY}. \just{divide por $\bY>0$}
\end{align*}
($\bY>0$ com probabilidade 1, pois cada $Y_i>0$.)

\subsubsection*{Passo 3 -- consistência}
\begin{enumerate}[label=(\roman*)]
\item Pela LGN, $\bY\convP\E(Y)=\frac1\lambda$.
\item A função $g(x)=\frac1x$ é contínua em todo $x\ne0$; em particular, em $x=\frac1\lambda>0$.
\item Pela aplicação contínua (Cont), $\tilde\lambda=g(\bY)\convP g\big(\frac1\lambda\big)=\frac{1}{1/\lambda}=\lambda$.
\end{enumerate}
Logo $\tilde\lambda$ é consistente.

\subsubsection*{Passo 4 -- viés}
Como $\tilde\lambda=\frac{1}{\bY}=\frac{n}{S}$, com $S=\sum_iY_i$, precisamos de $\E(1/S)$. Para isso precisamos da \textbf{distribuição de $S$}.

\textit{4a) FGM de uma $Y\sim\text{Exp}(\lambda)$.} Para $t<\lambda$:
\begin{align*}
M_Y(t)=\E(e^{tY})&=\int_0^\infty e^{ty}\lambda e^{-\lambda y}\,dy \just{P7}\\
&=\lambda\int_0^\infty e^{-(\lambda-t)y}\,dy \just{L3: $e^{ty}e^{-\lambda y}=e^{-(\lambda-t)y}$}\\
&=\lambda\cdot\frac{\Gamma(1)}{(\lambda-t)^1}=\frac{\lambda}{\lambda-t}. \just{I1 com $\alpha=1$; I2}
\end{align*}
A condição $t<\lambda$ garante que a taxa $\lambda-t$ é positiva, como I1 exige.

\textit{4b) FGM da soma.}
\begin{align*}
M_S(t)=\E\big(e^{t\sum_iY_i}\big)&=\E\Big(\prod_{i=1}^ne^{tY_i}\Big) \just{L3}\\
&=\prod_{i=1}^n\E(e^{tY_i}) \just{P6: independência}\\
&=\Big(\frac{\lambda}{\lambda-t}\Big)^n . \just{P8 e S5}
\end{align*}
A FGM de uma $\text{Gama}(\alpha,\beta)$ (taxa $\beta$) é $\big(\frac{\beta}{\beta-t}\big)^\alpha$. Pela unicidade (Pr4), $S\sim\text{Gama}(n,\lambda)$, ou seja,
\[
f_S(s)=\frac{\lambda^n}{\Gamma(n)}s^{n-1}e^{-\lambda s},\qquad s>0 .
\]

\textit{4c) Cálculo de $\E(1/S)$} (supondo $n\ge2$):
\begin{align*}
\E\Big(\frac1S\Big)&=\int_0^\infty\frac1s\cdot\frac{\lambda^n}{\Gamma(n)}s^{n-1}e^{-\lambda s}\,ds \just{P7}\\
&=\frac{\lambda^n}{\Gamma(n)}\int_0^\infty s^{n-2}e^{-\lambda s}\,ds \just{I3; $s^{-1}s^{n-1}=s^{n-2}$}\\
&=\frac{\lambda^n}{\Gamma(n)}\int_0^\infty s^{(n-1)-1}e^{-\lambda s}\,ds \just{reescreve no formato de I1}\\
&=\frac{\lambda^n}{\Gamma(n)}\cdot\frac{\Gamma(n-1)}{\lambda^{n-1}} \just{I1: $\alpha=n-1>0$}\\
&=\lambda^{n-(n-1)}\cdot\frac{\Gamma(n-1)}{\Gamma(n)} \just{$\lambda^a/\lambda^b=\lambda^{a-b}$}\\
&=\lambda\cdot\frac{\Gamma(n-1)}{(n-1)\,\Gamma(n-1)} \just{I2: $\Gamma(n)=(n-1)\Gamma(n-1)$}\\
&=\frac{\lambda}{n-1}.
\end{align*}

\textit{4d) Esperança e viés de $\tilde\lambda$.}
\begin{align*}
\E(\tilde\lambda)&=\E\Big(\frac nS\Big)=n\,\E\Big(\frac1S\Big) \just{P1}\\
&=\frac{n\lambda}{n-1},\\
\B(\tilde\lambda)&=\frac{n\lambda}{n-1}-\lambda=\frac{n\lambda-(n-1)\lambda}{n-1} \just{denominador comum}\\
&=\frac{n\lambda-n\lambda+\lambda}{n-1}=\frac{\lambda}{n-1}.
\end{align*}
Então $\tilde\lambda$ é \textbf{viciado} (superestima $\lambda$ em média), mas $\B(\tilde\lambda)=\frac{\lambda}{n-1}\to0$: é \textbf{assintoticamente não-viciado}.

\textit{Por que superestima?} A função $g(x)=1/x$ é convexa em $x>0$. Pela desigualdade de Jensen, $\E[g(\bY)]\ge g(\E\bY)$, isto é, $\E(1/\bY)\ge1/\E(\bY)=\lambda$. É a mesma lição da observação sobre funções não lineares: $\E(1/\bY)\neq1/\E(\bY)$.

\textit{Correção do viés.} Como $\E(\tilde\lambda)=\frac{n}{n-1}\lambda$, multiplicando por $\frac{n-1}{n}$ (P1):
\[
\E\Big(\frac{n-1}{n}\tilde\lambda\Big)=\frac{n-1}{n}\cdot\frac{n}{n-1}\lambda=\lambda .
\]
Logo $\tilde\lambda^*=\frac{n-1}{n}\cdot\frac nS=\frac{n-1}{S}$ é não-viciado.

\subsubsection*{Numérico}
Amostra $(0{,}4;\ 1{,}1;\ 0{,}7;\ 2{,}8)$, $n=4$:
\begin{align*}
\sum_iy_i&=0{,}4+1{,}1+0{,}7+2{,}8=5{,}0,\\
\bar y&=\frac{5{,}0}{4}=1{,}25,\\
\tilde\lambda&=\frac1{1{,}25}=0{,}8,\\
\tilde\lambda^*&=\frac{n-1}{\sum_iy_i}=\frac{3}{5{,}0}=0{,}6 .
\end{align*}

%---------------------------------------------------------------
\subsection{Exemplo 5 -- Momentos com dois parâmetros: Gama$(\alpha,\beta)$}\label{ex:gama2}

\textbf{Enunciado.} Seja $Y_1,\dots,Y_n$ a.a. de $Y\sim\text{Gama}(\alpha,\beta)$ (taxa), com $f(y)=\frac{\beta^\alpha}{\Gamma(\alpha)}y^{\alpha-1}e^{-\beta y}$ para $y>0$, e ambos os parâmetros desconhecidos. Encontre os estimadores de momentos.

\subsubsection*{Passo 1 -- dois parâmetros exigem dois momentos}
\textit{Primeiro momento.}
\begin{align*}
\E(Y)&=\int_0^\infty y\cdot\frac{\beta^\alpha}{\Gamma(\alpha)}y^{\alpha-1}e^{-\beta y}\,dy \just{P7}\\
&=\frac{\beta^\alpha}{\Gamma(\alpha)}\int_0^\infty y^{(\alpha+1)-1}e^{-\beta y}\,dy \just{I3; $y\cdot y^{\alpha-1}=y^{\alpha}=y^{(\alpha+1)-1}$}\\
&=\frac{\beta^\alpha}{\Gamma(\alpha)}\cdot\frac{\Gamma(\alpha+1)}{\beta^{\alpha+1}} \just{I1 com expoente $\alpha+1$}\\
&=\frac{\Gamma(\alpha+1)}{\Gamma(\alpha)}\cdot\frac{\beta^\alpha}{\beta^{\alpha+1}}
=\frac{\alpha\,\Gamma(\alpha)}{\Gamma(\alpha)}\cdot\frac1\beta \just{I2}\\
&=\frac\alpha\beta .
\end{align*}
\textit{Segundo momento.} Mesmo caminho, agora com $y^2\cdot y^{\alpha-1}=y^{(\alpha+2)-1}$:
\begin{align*}
\E(Y^2)&=\frac{\beta^\alpha}{\Gamma(\alpha)}\cdot\frac{\Gamma(\alpha+2)}{\beta^{\alpha+2}} \just{P7, I3, I1}\\
&=\frac{(\alpha+1)\,\alpha\,\Gamma(\alpha)}{\Gamma(\alpha)}\cdot\frac{1}{\beta^2} \just{I2 duas vezes}\\
&=\frac{\alpha(\alpha+1)}{\beta^2}.
\end{align*}
O passo ``I2 duas vezes'' é $\Gamma(\alpha+2)=(\alpha+1)\Gamma(\alpha+1)=(\alpha+1)\,\alpha\,\Gamma(\alpha)$.

\subsubsection*{Passo 2 -- montar o sistema}
\[
\text{(i)}\quad\frac{\tilde\alpha}{\tilde\beta}=M_1,
\qquad\qquad
\text{(ii)}\quad\frac{\tilde\alpha(\tilde\alpha+1)}{\tilde\beta^2}=M_2 .
\]

\subsubsection*{Passo 3 -- resolver}
\textit{Separe a fração de (ii) em duas:}
\[
\frac{\tilde\alpha(\tilde\alpha+1)}{\tilde\beta^2}=\frac{\tilde\alpha^2+\tilde\alpha}{\tilde\beta^2}=\frac{\tilde\alpha^2}{\tilde\beta^2}+\frac{\tilde\alpha}{\tilde\beta^2}.
\]
\textit{Use (i) ao quadrado:} $\frac{\tilde\alpha^2}{\tilde\beta^2}=\Big(\frac{\tilde\alpha}{\tilde\beta}\Big)^2=M_1^2$. Substituindo em (ii):
\begin{align*}
M_1^2+\frac{\tilde\alpha}{\tilde\beta^2}&=M_2\\
\text{(iii)}\qquad\frac{\tilde\alpha}{\tilde\beta^2}&=M_2-M_1^2 . \just{subtrai $M_1^2$}
\end{align*}
\textit{Divida (i) por (iii)}, membro a membro:
\begin{align*}
\frac{\tilde\alpha/\tilde\beta}{\tilde\alpha/\tilde\beta^2}&=\frac{M_1}{M_2-M_1^2}\\
\frac{\tilde\alpha}{\tilde\beta}\cdot\frac{\tilde\beta^2}{\tilde\alpha}&=\frac{M_1}{M_2-M_1^2} \just{inverso da fração}\\
\tilde\beta&=\frac{M_1}{M_2-M_1^2}. \just{cancela $\tilde\alpha$ e um $\tilde\beta$}
\end{align*}
\textit{Volte em (i)} para achar $\tilde\alpha$:
\[
\tilde\alpha=M_1\,\tilde\beta=M_1\cdot\frac{M_1}{M_2-M_1^2}=\frac{M_1^2}{M_2-M_1^2}.
\]

\textit{Interpretação do denominador.} Dividindo o Lema-SQ por $n$:
\[
\frac1n\sum_i(Y_i-\bY)^2=\frac1n\sum_iY_i^2-\bY^2=M_2-M_1^2 .
\]
É a variância amostral com divisor $n$: positiva, a menos que todas as observações sejam iguais.

\subsubsection*{Numérico}
Amostra $(1,2,2,3,7)$, $n=5$. Organize em tabela (é assim que se faz em prova):
\begin{center}
\begin{tabular}{@{}cccccc|c@{}}
\toprule
$i$ & 1 & 2 & 3 & 4 & 5 & Soma\\
\midrule
$y_i$ & 1 & 2 & 2 & 3 & 7 & 15\\
$y_i^2$ & 1 & 4 & 4 & 9 & 49 & 67\\
$(y_i-\bar y)^2$ & 4 & 1 & 1 & 0 & 16 & 22\\
\bottomrule
\end{tabular}
\end{center}
\begin{align*}
M_1&=\frac{15}{5}=3, \qquad M_2=\frac{67}{5}=13{,}4,\\
M_2-M_1^2&=13{,}4-9=4{,}4 \qquad\Big(\text{confere: }\tfrac{22}{5}=4{,}4\Big),\\
\tilde\beta&=\frac{3}{4{,}4}=\frac{15}{22}\approx0{,}682,\\
\tilde\alpha&=\frac{3^2}{4{,}4}=\frac{9}{4{,}4}=\frac{45}{22}\approx2{,}045 .
\end{align*}
\textit{Conferência} na equação (i): $\tilde\alpha/\tilde\beta=\frac{45/22}{15/22}=\frac{45}{15}=3=M_1$. Correto.

%---------------------------------------------------------------
\subsection{Exemplo 6 -- Momentos: Normal com $\mu$ e $\sigma^2$ desconhecidos}\label{ex:normalmm}

\textbf{Passo 1.} $\E(Y)=\mu$ e, por P4, $\E(Y^2)=\Var(Y)+[\E(Y)]^2=\sigma^2+\mu^2$.

\textbf{Passo 2 (sistema).}
\[
\text{(i)}\quad\tilde\mu=M_1,\qquad\qquad\text{(ii)}\quad\tilde\sigma^2+\tilde\mu^2=M_2 .
\]

\textbf{Passo 3 (resolver).} De (i), $\tilde\mu=\bY$. Substituindo em (ii):
\begin{align*}
\tilde\sigma^2&=M_2-\tilde\mu^2=M_2-M_1^2 \just{subtrai $\tilde\mu^2$}\\
&=\frac1n\sum_iY_i^2-\bY^2
=\frac1n\Big[\sum_iY_i^2-n\bY^2\Big] \just{põe $\frac1n$ em evidência}\\
&=\frac1n\sum_{i=1}^n(Y_i-\bY)^2 . \just{Lema-SQ}
\end{align*}

\textbf{Passo 4 (esperança).}
\begin{align*}
\E(\tilde\sigma^2)&=\frac1n\,\E\Big[\sum_i(Y_i-\bY)^2\Big] \just{P1}\\
&=\frac1n(n-1)\sigma^2=\frac{n-1}{n}\sigma^2 . \just{Teo-SQ}
\end{align*}
O estimador de momentos da variância usa divisor $n$, e não $n-1$. A sua variância e o seu EQM ficam para a lista.

%---------------------------------------------------------------
\subsection{Exemplo 7 -- Momentos: quando a estimativa contradiz os dados}\label{ex:unifdesloc}

\textbf{Enunciado.} Seja $Y\sim U(\theta,\theta+1)$, $\theta\in\R$. A densidade é $f(y;\theta)=\frac{1}{(\theta+1)-\theta}=1$ para $\theta<y<\theta+1$, e zero fora desse intervalo.

\subsubsection*{Passo 1 -- momentos}
\begin{align*}
\E(Y)&=\int_\theta^{\theta+1}y\cdot1\,dy=\Big[\frac{y^2}{2}\Big]_\theta^{\theta+1} \just{P7, I4}\\
&=\frac{(\theta+1)^2-\theta^2}{2}
=\frac{\theta^2+2\theta+1-\theta^2}{2} \just{$(a+b)^2=a^2+2ab+b^2$}\\
&=\frac{2\theta+1}{2}=\theta+\frac12 .
\end{align*}
\begin{align*}
\E(Y^2)&=\Big[\frac{y^3}{3}\Big]_\theta^{\theta+1}=\frac{(\theta+1)^3-\theta^3}{3} \just{P7, I4}\\
&=\frac{\theta^3+3\theta^2+3\theta+1-\theta^3}{3} \just{$(a+b)^3=a^3+3a^2b+3ab^2+b^3$}\\
&=\theta^2+\theta+\frac13,
\end{align*}
\begin{align*}
\Var(Y)&=\E(Y^2)-[\E(Y)]^2 \just{P4}\\
&=\theta^2+\theta+\frac13-\Big(\theta^2+\theta+\frac14\Big)
=\frac13-\frac14=\frac1{12}.
\end{align*}

\subsubsection*{Passo 2 -- estimador de momentos}
\[
\tilde\theta+\frac12=\bY\quad\Longrightarrow\quad\tilde\theta=\bY-\frac12 .
\]

\subsubsection*{Passo 3 -- propriedades}
\begin{align*}
\E(\tilde\theta)&=\E(\bY)-\frac12 \just{P1}\\
&=\theta+\frac12-\frac12=\theta, \just{Teorema~\ref{teo:media}}\\
\Var(\tilde\theta)&=\Var(\bY)=\frac{1/12}{n}=\frac{1}{12n}. \just{P2 com $a=1$; Teorema~\ref{teo:media}}
\end{align*}
Então $\tilde\theta$ é não-viciado, com $\EQM=\frac1{12n}\to0$: é consistente.

\subsubsection*{Passo 4 -- o problema com os dados}
Amostra $(0{,}9;\ 1{,}0;\ 1{,}1;\ 1{,}85)$:
\[
\bar y=\frac{0{,}9+1{,}0+1{,}1+1{,}85}{4}=\frac{4{,}85}{4}=1{,}2125,
\qquad
\tilde\theta=1{,}2125-0{,}5=0{,}7125 .
\]
\textit{Quais valores de $\theta$ são compatíveis com a amostra?} Cada observação precisa cair dentro do suporte, $\theta<y_i<\theta+1$. Separamos as duas desigualdades:
\begin{itemize}
\item $\theta<y_i$ para \textbf{todo} $i$ $\iff$ $\theta<\min_iy_i=0{,}9$;
\item $y_i<\theta+1$ para \textbf{todo} $i$ $\iff$ $\theta>y_i-1$ para todo $i$ $\iff$ $\theta>\max_iy_i-1=1{,}85-1=0{,}85$.
\end{itemize}
Só valores de $\theta$ em $(0{,}85;\ 0{,}9)$ são compatíveis com os dados. A verossimilhança mostra o mesmo:
\[
L(\theta)=\prod_{i=1}^n\ind\{\theta<y_i<\theta+1\}=\ind\{\max_iy_i-1<\theta<\min_iy_i\}.
\]
Ela vale 1 nesse intervalo e 0 fora dele.

Como $\tilde\theta=0{,}7125<0{,}85$, temos $\tilde\theta+1=1{,}7125<1{,}85$. Ou seja, o modelo ajustado diria que a observação $1{,}85$ era \textbf{impossível}: $f(1{,}85;\tilde\theta)=0$ e $L(\tilde\theta)=0$.

\textbf{Lição:} o método dos momentos só olha para $\bY$ e ignora a informação dos extremos da amostra. Quando o suporte depende do parâmetro, sempre confira a compatibilidade e compare com estimadores baseados em $Y_{(1)}$ ou $Y_{(n)}$.

%=====================================================================
\section{Roteiros de prova (checklist)}
%=====================================================================

\begin{mdframed}[linecolor=Gray,backgroundcolor=Gray!5]
\textbf{Mostrar que $T$ é não-viciado.}
Calcule $\E_\theta(T)$ usando linearidade (escreva $T$ como combinação de $Y_i$ sempre que possível) e mostre que é igual a $\theta$ para todo $\theta$.

\medskip
\textbf{Encontrar $\Var(T)$.}
Se $T=\sum a_iY_i$ com $Y_i$ independentes: $\sum a_i^2\Var(Y_i)$. Se $T$ envolve somas de quadrados sob normalidade: escreva $T$ como constante $\times\ \chi^2$. Se $T$ é máximo ou mínimo: obtenha a densidade pelo Teorema~\ref{teo:ordem} e integre para $\E(T)$ e $\E(T^2)$.

\medskip
\textbf{Encontrar o EQM.}
$\EQM=\Var+\B^2$. Nunca esqueça o viés ao quadrado.

\medskip
\textbf{Consistência em MQ.}
Mostre que $\lim_{n\to\infty}\EQM_\theta(T_n)=0$ para cada $\theta$ fixo.

\medskip
\textbf{Consistência em probabilidade.}
Ou (i) cite MQ $\Rightarrow$ P, escrevendo a cadeia de Chebyshev
$\Prob(|T_n-\theta|\ge\varepsilon)\le\EQM/\varepsilon^2\to0$;
ou (ii) calcule $\Prob(|T_n-\theta|\ge\varepsilon)$ diretamente pela f.d.a.;
ou (iii) use LGN + aplicação contínua.

\medskip
\textbf{Não-consistência.}
Exiba um $\varepsilon>0$ com $\Prob(|T_n-\theta|\ge\varepsilon)\not\to0$.

\medskip
\textbf{Comparar estimadores.}
Calcule os dois EQM, forme a diferença ou a razão e estude o sinal \emph{como função de $\theta$} e de $n$. Conclua se há dominância uniforme ou cruzamento, e discuta também viés e consistência.

\medskip
\textbf{Estimador de momentos.}
Roteiro da seção do método dos momentos: momentos populacionais $\to$ igualar aos amostrais $\to$ resolver $\to$ checar $\Theta$ e compatibilidade com os dados.
\end{mdframed}

\end{document}